\documentclass{handout}

\assignment{Lab 6}
\title{Ideal gas law}

\begin{document}
\maketitle

Go to the Phet simulation ``gas properties''.

\section{Pressure and amount}

\begin{questions}
    \question Add light molecules in units of 50 to the container. Let the pressure gauge adjust and make note of the pressure. Plot several data points (4 ish) until you see a trend.

    \begin{figure}[h!]
        \includegraphics[width=0.5\linewidth]{Images/PN_diagram.pdf}    \end{figure}

    \question Draw a trendline and find the equation of that line: $P=aN+b$, where $a,b$ are the slope and y-intercept, respectively. What are the values and units of the trendline parameters?

    \answerspace{1.5in}
    \clearpage

    \question Explain the relationship between pressure and number of molecules in a gas (when volume and temperature are held constant) in your own words.

    \answerspace{1.5in}

    \question Try switching the light for heavy particles and see what happens to the pressure. Is ideal gas behavior primarily dependent on the total mass of the gas or simply the number of molecules it contains? Explain your reasoning.

    \answerspace{1.5in}
\end{questions}

\section{Pressure and temperature}
Put 500 light molecules in the container.
\begin{questions}
    \question Set the thermometer to measure in $^\circ$ C. Use the heat/cool option at the bottom of the screen to explore the pressure in a neighborhood around human temperatures (from -20 to 40 C).
    \begin{figure}[h!]
        \includegraphics[width=0.5\linewidth]{Images/PT_diagram.pdf}    \end{figure}
    \clearpage

    \question Again, draw the trendline and find the equation of the line: $P=aT+b$.

    \answerspace{1.5in}

    \question From your equation, find the temperature at which the gas would exert no force on the walls of the container. Express your answer in $^\circ$ C.

    \answerspace{1.5in}

    \question In your own words, explain the relationship between pressure and temperature. Explain also why it is more useful to use the Kelvin scale rather than Celsius for cases such as this.

    \answerspace{1.5in}

\end{questions}

\clearpage
\section{Types of processes in gases}
Reset the simulation (orange button) and put 250 light molecules in the container. We will be measuring the pressure and volume, so tick the box to show the width of the container. The cross-sectional area of the container won't be changing, so the volume is $V=Aw$ proportional to the width. We will plot width as a stand-in for volume.

\begin{questions}
\question Use the handle at the left side of the screen to change the volume of the box. You can control this between 5 and 15 nm, so consider a few points in that range and plot the pressures that result.

    \begin{figure}[h!]
        \includegraphics[width=0.5\linewidth]{Images/PV_diagram.pdf}    \end{figure}

\question Cool the system to around 150 K and repeat the previous measurements, plotting them on the same diagram.

\answerspace{1.5in}

\question Using the ideal gas law, explain how you would read the temperature of a gas from its location on a PV diagram as above. \label{PVreadT}

\answerspace{1.5in}
\clearpage

\question In which case, compression or expansion of the gas, is \textbf{work} ($\vec{\boldsymbol{F}}\cdot \boldsymbol{\Delta \vec{r}}$) being done \textbf{on} the gas?

\answerspace{1in}

\question If the temperature is held constant, the thermal energy in the gas does not change. If a gas expands at constant temperature, does heat energy flow into or out of the gas?

\answerspace{1in}

\question Set the gas to be held at constant \textbf{pressure} (the option with T arrows). Heat and cool the gas at constant pressure. Draw the version on the graph below where the gas is being heated (use arrows on your graph to denote before and after).

\begin{figure}[h!]
        \includegraphics[width=0.5\linewidth]{Images/PV_diagram.pdf}    \end{figure}

\question Does this fit with your answer to \ref{PVreadT}? Why or why not?

\clearpage
\end{questions}

Use the following graph to answer the remaining questions
\begin{figure}
    \centering
    \includegraphics[width=0.5\linewidth]{Images/CycleDiagram.pdf}\end{figure}

\begin{questions}
    \question In process A, heat energy flows \hspace{.3in}into\hspace{.3in}/\hspace{.3in} out of \hspace{.3in} the gas at constant \blank[1.5in] \answerspace{\baselineskip}
    \question In process B, heat energy flows \hspace{.3in}into\hspace{.3in}/\hspace{.3in} out of \hspace{.3in} the gas at constant \blank[1.5in] \answerspace{\baselineskip}
    \question In process C, heat energy flows \hspace{.3in}into\hspace{.3in}/\hspace{.3in} out of \hspace{.3in} the gas at constant \blank[1.5in] \answerspace{\baselineskip}
    \question In process D, heat energy flows \hspace{.3in}into\hspace{.3in}/\hspace{.3in} out of \hspace{.3in} the gas at constant \blank[1.5in] \answerspace{\baselineskip}
    \question \hspace{.3in}T\hspace{.3in}/\hspace{.3in} F \hspace{.3in}: The gas is in exactly the same state before and after going through all four processes.
\end{questions}

\end{document}
